\textbf{Guaranteed intervals.} If $a\in[L,U]$ and $b\in[V,W]$, then $a+b\in[L+V,U+W]$ and $a-b\in[L-W,U-V]$. When a and b may vary freely in their separate closed intervals, these bounds are tight and both endpoints can occur. This is deterministic uncertainty: no probability distribution or confidence percentage is assumed.

``Always reaches r'' means the minimum is at least r; ``always fits within s'' means the maximum is at most s. Equality is allowed in both. Except when a problem explicitly asks for whole-unit pairs, hidden endpoints may lie anywhere in the band, including between ticks and at the edges. The cards stipulate guaranteed bounds; the ruler itself does not establish an instrument error bound.

\textbf{Dependencies matter.} Repeated A is one unchanged quantity. Thus (A+4)-(A+1)=3 and (A+B)-A=B. Using separate ranges for repeated copies can enclose the answer too broadly. A known relation such as A+B=12 restricts possible endpoint combinations. The elementary formulas remain enclosures, but need not remain tight when the variables are linked.

\textbf{Discovery by band.} K-1 builds extreme examples and tries to defeat ``always'' claims. Grades 2-3 designs the widest safe card and compares shared versus separate strips. Grades 4-5 tightens bounds using an exact total and shared-piece cancellation. Testing examples can disprove a claim; endpoint arguments prove a guarantee and its sharpness.
